% ============================================================
% M3 S3 练习件·W2 臂A 片件 —— 课时06 两条直线的位置关系（2.2.3）
% 生成：M3 成卷轮 S3 W2 臂A｜2026-09-14｜编译 cwd＝本目录（中文路径禁 \input 绝对路径）
%   xelatex main-true.tex ×2 → 含详解印本；xelatex main-false.tex ×2 → 纯题版（单源双本）
% 输入链（全只读）：题面库 课时06-两条直线的位置关系.md（题面逐字装配）＋ -答案侧.md（值逐字回嵌）
%   ＋ 定稿/批B-课时06-两条直线的位置关系.md（详解回嵌源）；toolchain＝qp-m3.sty（钉后版
%   md5 7c3930362be8a0a2bdf21bbf8ac16573，本地挂载副本，破损即停）。
% 母版体例：整目录照抄 成卷/练习件/课时01-坐标法/（骨架/宏用法/门谱原样，只换内容层）。
% 键账：件内 ansblock 键集 ≡ 题面库片练键集 ≡ manifest 练键序 T1~T16（16 键，零缺漏零多余；
%   本片键式＝T 式，承 canonical 键名总表批B 口径）；导学键 G1~G5 不入本件（导学件域）；
%   拓区隔离：2章-拓/测/滚 键不入本件（S4 拓展册收）。
% 卷式例外案B：练习件不属卷式——答案块物理落点恒为题后 ansblock，禁卷末附卷制；
%   全线括线模（导言 \ansblockgrayfalse 一次置定，件内不切换；灰底盒不可跨栏断，练习件禁用）。
% 题序＝manifest 练键序 T1~T16（零重排）；印面题号 1~16 连号（\tihao 号＝\ansitem 号同键）；
%   三档切位 1~7／8~14／15~16（照库序切位，非难度重排）。
% 槽配实录：单选8（T1~T7、T9）＋多选2（T8、T10，≤4 合规）＋填空4（T11~T14）＋解答2（T15~T16）；
%   波次方案基线 10选4填2解 为槽配基线，本片实配照题面侧头标（批B §二 配额对账同口径）。
% 替换题注记：T3/T4 系双收补换位题（2章件4-#9／2章件2-#36，0913 落盘），题面侧【注】行
%   （换位留痕）不入印面；T14 白名单注同不入。
% 需选学注：T13/T16 系§2.2.4(06B) 前引注题，题面侧「前置提示」行照批B 加注口径印入题面
%   （05 片 E13/E14 预习注同式：注行括注置 \tieside 后、题面前；与导学件 06B 件首 \zhuzhu
%   「课时06 件内 T13/T16 已加前引注」承接口径一致）。
% 选项槽位：默认四连排（0.25\linewidth−0.25\lxhang），超宽降档梯 四连排→两连排
%   （0.5\linewidth−0.5\lxhang−0.25em）→单列 \lxopt；禁缩字号/负 kern/删标点凑宽。
%   本片实测降档：题1/4/5/6 两连排、题8/10 单列，余四连排（真算读数见值台账「降档登记」）。
% 解答书写区：T15/T16 均两问 \liubai[24mm]（默认 16mm、两问 24mm、三问 32mm、cap 40mm；
%   纯题档吞 ansblock 不吞 \liubai）。
% 填空印答：T11~T14 题面空位 \kongda{值}（详解版印答、纯题版退定宽空线，两档等形）；
%   多空位逐空 \kongda，值取答案侧「值：」行同串按空位拆分（印答与 ansitem 同源零漂）。
% 难度词口径：题面侧头标第 4 字段未带档位词者，照题面库全库主档映射 0.94→简单、
%   0.85→中档、0.65→中档（带档位词者照录；登记见值台账「难度词口径」）。
% 知识点预排注记：导学件课时06 尚未落盘（目录仅 sty），知识点 N 系预排——
%   一＝两直线平行的判定、二＝两直线垂直的判定、三＝交点与位置关系综合应用，
%   照批B 课时切分与题型名主干分派；导学件落盘后如异动按变更协议回核。
% ============================================================
\documentclass[fontset=none]{ctexart}
\usepackage{qp-m3}
% —— 印面逐字符号路由补钉（承导学件母版；− 走 TNR、▱ NSC 子块；禁改 sty 保 md5） ——
\xeCJKDeclareCharClass{Default}{"2212}%
\setCJKfamilyfont{nscblk}[Path=C:/提示词/工作区/字替对照-0909/variantF/fonts/,BoldFont=NSC-w600.ttf]{NSC-w500.ttf}%
\xeCJKDeclareSubCJKBlock{nscblk}{"25B1}%
\newcommand{\pxparallelogram}{{\CJKfamily{nscblk}▱}}%
% —— 断行微调（承导学件母版实测档，三零门承重；\emergencystretch 不另设） ——
\xeCJKsetup{CJKglue={\hskip 0.10em plus 0.08em minus 0.05em}}%
% —— 练习件全线括线模（S3 波次方案 §四.3：导言一次置定、件内不切换） ——
\ansblockgrayfalse
% —— 页脚件名（qp-headfoot 复用入口） ——
\renewcommand{\qpzhangming}{第二章\quad 平面解析几何}%
\renewcommand{\qpjianming}{练习件}%

\begin{document}

\keshi{第6课时\quad 两条直线的位置关系}
\begin{multicols}{2}
\emergencystretch=1em

% ============================================================
% 基础巩固（题1~7＝T1~T7：单选7；题后紧跟答案制，全线括线 ansblock）
% ============================================================
\dangbadge{基}{础}{巩}{固}

\tihao{1}\zu{两直线位置关系的判定×单选} \tieside{中档(知识点二)}直线\(l_1\)：\(2ax+y-2=0\)与直线\(l_2\)：\(x-(a+1)y+2=0\)互相垂直，则这两条直线的交点坐标为\nobreak(\kongwei)
\bindopt {\fontsize{10.5pt}{\qpTJgLeadLiti}\selectfont \makebox[\dimexpr0.5\linewidth-0.5\lxhang-0.25em\relax][l]{A．\((-2/5,-6/5)\)}\makebox[\dimexpr0.5\linewidth-0.5\lxhang-0.25em\relax][l]{B．\((2/5,6/5)\)}\par}
\bindopt {\fontsize{10.5pt}{\qpTJgLeadLiti}\selectfont \makebox[\dimexpr0.5\linewidth-0.5\lxhang-0.25em\relax][l]{C．\((2/5,-6/5)\)}\makebox[\dimexpr0.5\linewidth-0.5\lxhang-0.25em\relax][l]{D．\((-2/5,6/5)\)}\par}

\begin{ansblock}[2章-练-课时06-T1]
% ans:2章-练-课时06-T1
\ansitem{1}{B}
\ansnote{详解}{\(A_1A_2+B_1B_2=2a\cdot 1+1\cdot (-(a+1))=a-1=0\)，得\(a=1\)（\(a=-1\)时\(l_2\)：\(x+2=0\)铅直而\(l_1\)斜率2，不垂直，已排除）。此时\(l_1\)：\(2x+y-2=0\)，\(l_2\)：\(x-2y+2=0\)。由\(y=2-2x\)代入：\(x-2(2-2x)+2=0\)，\(5x=2\)，\(x=2/5\)，\(y=6/5\)。交点\((2/5,6/5)\)，选B。}
\end{ansblock}

\tihao{2}\zu{两直线位置关系的判定×单选} \tieside{中档(知识点一)}满足下列条件的直线\(l_1\)与\(l_2\)，其中\(l_1\parallel l_2\)的是\nobreak(\kongwei)
\bindp ①\(l_1\)的斜率为2，\(l_2\)过点\(A(1,2)\)、\(B(4,8)\)；
\bindp ②\(l_1\)经过点\(P(3,3)\)、\(Q(-5,3)\)，\(l_2\)平行于\(x\)轴，但不经过\(P\)点；
\bindp ③\(l_1\)经过点\(M(-1,0)\)、\(N(-5,-2)\)，\(l_2\)经过点\(R(-4,3)\)、\(S(0,5)\)．
\bindopt {\fontsize{10.5pt}{\qpTJgLeadLiti}\selectfont \makebox[\dimexpr0.5\linewidth-0.5\lxhang-0.25em\relax][l]{A．①②}\makebox[\dimexpr0.5\linewidth-0.5\lxhang-0.25em\relax][l]{B．②③}\par}
\bindopt {\fontsize{10.5pt}{\qpTJgLeadLiti}\selectfont \makebox[\dimexpr0.5\linewidth-0.5\lxhang-0.25em\relax][l]{C．①③}\makebox[\dimexpr0.5\linewidth-0.5\lxhang-0.25em\relax][l]{D．①②③}\par}

\begin{ansblock}[2章-练-课时06-T2]
% ans:2章-练-课时06-T2
\ansitem{2}{B}
\ansnote{详解}{①\(k_{AB}=(8-2)/(4-1)=2=k_1\)，但只知斜率相等，\(l_1\)与\(l_2\)可能重合，①不入选。②\(k_{PQ}=0\)，\(l_1\)为水平线；\(l_2\)平行于\(x\)轴且不过\(P\)，两水平线不重合，平行，②入选。③\(k_1=(-2-0)/(-5+1)=1/2\)，\(k_2=(5-3)/(0+4)=1/2\)；检验\(M(-1,0)\)不在\(l_2\)（\(l_2\)：\(y=x/2+5\)，\(x=-1\)时\(y=9/2\)）上，故平行，③入选。选B。}
\end{ansblock}

\tihao{3}\zu{两直线位置关系的判定×单选} \tieside{中档(知识点一)}已知直线\(l_1\)过\((0,0)\)、\((1,-3)\)两点，直线\(l_2\)的方程为\(ax+y-2=0\)，如果\(l_1\parallel l_2\)，则\(a\)值为\nobreak(\kongwei)
\bindopt {\fontsize{10.5pt}{\qpTJgLeadLiti}\selectfont \makebox[\dimexpr0.25\linewidth-0.25\lxhang\relax][l]{A．\(-3\)}\makebox[\dimexpr0.25\linewidth-0.25\lxhang\relax][l]{B．\(1/3\)}\makebox[\dimexpr0.25\linewidth-0.25\lxhang\relax][l]{C．\(-1/3\)}\makebox[\dimexpr0.25\linewidth-0.25\lxhang\relax][l]{D．\(3\)}\par}

\begin{ansblock}[2章-练-课时06-T3]
% ans:2章-练-课时06-T3
\ansitem{3}{D}
\ansnote{详解}{\(k_1=(-3-0)/(1-0)=-3\)；\(l_2\)：\(ax+y-2=0\)的斜率\(k_2=-a\)。\(l_1\parallel l_2\Longleftrightarrow k_1=k_2\Longrightarrow -a=-3\Longrightarrow a=3\)。重合检验：\(a=3\)时\(l_2\)：\(3x+y-2=0\)，取\(l_1\)上点\((0,0)\)代入得\(-2\neq 0\)，\((0,0)\)不在\(l_2\)上，两直线不重合，平行成立。选D。}
\end{ansblock}

\tihao{4}\zu{两直线位置关系的判定×单选} \tieside{中档(知识点二)}已知直线\(l_1\)的倾斜角为\(60^\circ\)，直线\(l_2\)经过点\(A(1,\sqrt{3})\)、\(B(-2,2\sqrt{3})\)，则直线\(l_1\)与\(l_2\)的位置关系是\nobreak(\kongwei)\par
\bindopt {\fontsize{10.5pt}{\qpTJgLeadLiti}\selectfont \makebox[\dimexpr0.5\linewidth-0.5\lxhang-0.25em\relax][l]{A．平行或重合}\makebox[\dimexpr0.5\linewidth-0.5\lxhang-0.25em\relax][l]{B．平行}\par}
\bindopt {\fontsize{10.5pt}{\qpTJgLeadLiti}\selectfont \makebox[\dimexpr0.5\linewidth-0.5\lxhang-0.25em\relax][l]{C．垂直}\makebox[\dimexpr0.5\linewidth-0.5\lxhang-0.25em\relax][l]{D．重合}\par}

\begin{ansblock}[2章-练-课时06-T4]
% ans:2章-练-课时06-T4
\ansitem{4}{C}
\ansnote{详解}{\(k_1=\tan 60^\circ=\sqrt{3}\)；\(k_2=(2\sqrt{3}-\sqrt{3})/(-2-1)=\sqrt{3}/(-3)=-\sqrt{3}/3\)。\(k_1\cdot k_2=\sqrt{3}\times(-\sqrt{3}/3)=-1\)，且\(k_1\neq k_2\)（两线不平行、更不重合），故\(l_1\perp l_2\)。选C。}
\end{ansblock}

\tihao{5}\zu{两直线位置关系的判定×单选} \tieside{中档(知识点二)}已知\(\triangle ABC\)的顶点\(B(2,1)\)，\(C(-6,3)\)，其垂心为\(H(-3,2)\)，则其顶点\(A\)的坐标为\nobreak(\kongwei)
\bindopt {\fontsize{10.5pt}{\qpTJgLeadLiti}\selectfont \makebox[\dimexpr0.5\linewidth-0.5\lxhang-0.25em\relax][l]{A．\((-19,-62)\)}\makebox[\dimexpr0.5\linewidth-0.5\lxhang-0.25em\relax][l]{B．\((19,-62)\)}\par}
\bindopt {\fontsize{10.5pt}{\qpTJgLeadLiti}\selectfont \makebox[\dimexpr0.5\linewidth-0.5\lxhang-0.25em\relax][l]{C．\((-19,62)\)}\makebox[\dimexpr0.5\linewidth-0.5\lxhang-0.25em\relax][l]{D．\((19,62)\)}\par}

\begin{ansblock}[2章-练-课时06-T5]
% ans:2章-练-课时06-T5
\ansitem{5}{A}
\ansnote{详解}{\(k_{BC}=2/(-8)=-1/4\Longrightarrow k_{AH}=4\)；\par\noindent \(k_{BH}=1/(-5)=-1/5\Longrightarrow k_{AC}=5\)。\par\noindent 设\(A(x,y)\)：\((y-2)/(x+3)=4\)且\((y-3)/(x+6)=5\)\par\noindent \(\Longrightarrow y=4x+14\)且\(y=5x+33\Longrightarrow x=-19\)，\(y=-62\)。\par\noindent 验第三组垂直：\(k_{CH}=(2-3)/(-3+6)=-1/3\)，而\(k_{AB}=(1+62)/(2+19)=3\)，乘积为\(-1\)。选A。}
\end{ansblock}

\tihao{6}\zu{两直线位置关系的判定×单选} \tieside{中档(知识点二)}若直线\(kx+(1-k)y-3=0\)和直线\((k-1)x+(2k+3)y-2=0\)互相垂直，则\(k=\)\nobreak(\kongwei)
\bindopt {\fontsize{10.5pt}{\qpTJgLeadLiti}\selectfont \makebox[\dimexpr0.5\linewidth-0.5\lxhang-0.25em\relax][l]{A．\(-3\)或\(-1\)}\makebox[\dimexpr0.5\linewidth-0.5\lxhang-0.25em\relax][l]{B．\(3\)或\(1\)}\par}
\bindopt {\fontsize{10.5pt}{\qpTJgLeadLiti}\selectfont \makebox[\dimexpr0.5\linewidth-0.5\lxhang-0.25em\relax][l]{C．\(-3\)或\(1\)}\makebox[\dimexpr0.5\linewidth-0.5\lxhang-0.25em\relax][l]{D．\(-1\)或\(3\)}\par}

\begin{ansblock}[2章-练-课时06-T6]
% ans:2章-练-课时06-T6
\ansitem{6}{C}
\ansnote{详解}{垂直\(\Longleftrightarrow k(k-1)+(1-k)(2k+3)=0\)。展开：\(k^{2}-k+2k+3-2k^{2}-3k=-k^{2}-2k+3=0\)，即\(k^{2}+2k-3=0\)，\(k=1\)或\(k=-3\)。选C。}
\end{ansblock}

\tihao{7}\zu{两直线位置关系的判定×单选} \tieside{中档(知识点二)}已知\(a\neq 0\)，直线\(ax+(b+2)y+4=0\)与直线\(ax+(b-2)y-3=0\)互相垂直，则\(ab\)的最大值为\nobreak(\kongwei)
\bindopt {\fontsize{10.5pt}{\qpTJgLeadLiti}\selectfont \makebox[\dimexpr0.25\linewidth-0.25\lxhang\relax][l]{A．\(0\)}\makebox[\dimexpr0.25\linewidth-0.25\lxhang\relax][l]{B．\(2\)}\makebox[\dimexpr0.25\linewidth-0.25\lxhang\relax][l]{C．\(4\)}\makebox[\dimexpr0.25\linewidth-0.25\lxhang\relax][l]{D．\(\sqrt{2}\)}\par}

\begin{ansblock}[2章-练-课时06-T7]
% ans:2章-练-课时06-T7
\ansitem{7}{B}
\ansnote{详解}{垂直\(\Longleftrightarrow a\cdot a+(b+2)(b-2)=0\)，即\(a^{2}+b^{2}=4\)。由\(a^{2}+b^{2}\geqslant 2|ab|\)得\(|ab|\leqslant 2\)；取\(a=b=\sqrt{2}\)（满足\(a^{2}+b^{2}=4\)，系数检验：两直线分别为\(\sqrt{2}x+(2+\sqrt{2})y+4=0\)与\(\sqrt{2}x+(\sqrt{2}-2)y-3=0\)，\(A_1A_2+B_1B_2=2+(2+\sqrt{2})(\sqrt{2}-2)=2+(2-4)=0\)，垂直成立）时\(ab=2\)。故\(ab\)的最大值为2，选B。}
\end{ansblock}

% ============================================================
% 综合提升（题8~14＝T8~T14：多选2＋单选1＋填空4；题后紧跟答案制，全线括线 ansblock）
% ============================================================
\dangbadge{综}{合}{提}{升}

\tihao{8}\duoxuan\hspace{0.6em}\zu{两直线位置关系的判定×多选} \tieside{中档(知识点三)}已知直线\(l_1\)：\(x+ay-a=0\)和直线\(l_2\)：\(ax-(2a-3)y-1=0\)，下列说法正确的是\nobreak(\kongwei)

\lxopt{A．\(l_2\)始终过定点\((2/3,1/3)\)}

\lxopt{B．若\(l_1\parallel l_2\)，则\(a=1\)或\(-3\)}

\lxopt{C．若\(l_1\perp l_2\)，则\(a=0\)或\(2\)}

\lxopt{D．当\(a>0\)时，\(l_1\)始终不过第三象限}

\begin{ansblock}[2章-练-课时06-T8]
% ans:2章-练-课时06-T8
\ansitem{8}{ACD}
\ansnote{详解}{A：\(l_2\)按\(a\)整理：\(a(x-2y)+(3y-1)=0\)，令\(x-2y=0\)且\(3y-1=0\)得\((2/3,1/3)\)，恒过，A对。\par\noindent B：\(a=1\)时\(l_1\)：\(x+y-1=0\)，\(l_2\)：\(x+y-1=0\)，重合非平行；\(a=-3\)时\(k_1=-1/3\)，\(l_2\)：\(-3x+9y-1=0\)的斜率为\(1/3\)，不平行，B错。\par\noindent C：\(A_1A_2+B_1B_2=a+a(3-2a)=2a(2-a)=0\Longrightarrow a=0\)或2，C对。\par\noindent D：\(a>0\)时\(l_1\)：\(y=-x/a+1\)，斜率为负、纵截距\(1>0\)、横截距\(a>0\)，图象过一、二、四象限，不过第三象限，D对。选ACD。}
\end{ansblock}

\tihao{9}\zu{两直线位置关系的判定×单选} \tieside{中档(知识点三)}过直线\(l_1\)：\(2x+y-3=0\)与\(l_2\)：\(x-3y+2=0\)的交点，并与\(l_1\)垂直的直线的方程为\nobreak(\kongwei)
\bindopt {\fontsize{10.5pt}{\qpTJgLeadLiti}\selectfont \makebox[\dimexpr0.5\linewidth-0.5\lxhang-0.25em\relax][l]{A．\(x-2y-1=0\)}\makebox[\dimexpr0.5\linewidth-0.5\lxhang-0.25em\relax][l]{B．\(x-2y+1=0\)}\par}
\bindopt {\fontsize{10.5pt}{\qpTJgLeadLiti}\selectfont \makebox[\dimexpr0.5\linewidth-0.5\lxhang-0.25em\relax][l]{C．\(x+2y-1=0\)}\makebox[\dimexpr0.5\linewidth-0.5\lxhang-0.25em\relax][l]{D．\(x+2y+1=0\)}\par}

\begin{ansblock}[2章-练-课时06-T9]
% ans:2章-练-课时06-T9
\ansitem{9}{B}
\ansnote{详解}{联立\(2x+y-3=0\)，\(x-3y+2=0\)：由第一式\(y=3-2x\)代入：\(x-9+6x+2=0\)，\(x=1\)，\(y=1\)，交点\((1,1)\)。\(k_1=-2\)，所求斜率\(k=1/2\)。\(y-1=(1/2)(x-1)\)，即\(x-2y+1=0\)，选B。}
\end{ansblock}

\tihao{10}\duoxuan\hspace{0.6em}\zu{两直线位置关系的判定×多选} \tieside{中档(知识点三)}瑞士数学家欧拉1765年提出：任意三角形的外心、重心、垂心在同一条直线上，后人称这条直线为欧拉线。若已知\(\triangle ABC\)的顶点\(A(-4,0)\)，\(B(0,4)\)，其欧拉线方程为\(x-y+2=0\)，则下列正确的是\nobreak(\kongwei)

\lxopt{A．\(\triangle ABC\)重心的坐标为\((-1/3,2/3)\)或\((-2/3,1/3)\)}

\lxopt{B．\(\triangle ABC\)垂心的坐标为\((0,2)\)或\((-2,0)\)}

\lxopt{C．\(\triangle ABC\)顶点\(C\)的坐标为\((2,0)\)或\((0,-2)\)}

\lxopt{D．欧拉线将\(\triangle ABC\)分成的两部分的面积之比为\(4/5\)}

\begin{ansblock}[2章-练-课时06-T10]
% ans:2章-练-课时06-T10
\ansitem{10}{BCD}
\ansnote{详解}{\(AB\)中点\((-2,2)\)，\(k_{AB}=1\)，中垂线\(y-2=-(x+2)\)，即\(x+y=0\)。与欧拉线联立：\(x=-1\)，\(y=1\)，外心\(O'(-1,1)\)，\(R^{2}=9+1=10\)。\par\noindent 设\(C(m,n)\)：重心\(((m-4)/3,(n+4)/3)\)在欧拉线上\(\Longrightarrow m-n-2=0\)；\(|O'C|^{2}=10\Longrightarrow (m+1)^{2}+(n-1)^{2}=10\)，即\(m^{2}+n^{2}+2m-2n=8\)。代入\(n=m-2\)：\(2m^{2}-4m=0\Longrightarrow m=0\)或2，\(C(2,0)\)或\((0,-2)\)，C对。\par\noindent 重心为\((-2/3,4/3)\)或\((-4/3,2/3)\)，A给的数值错，A错。\par\noindent 垂心：由欧拉关系\(H=3G-2O'\)（\(G\)为重心）：\(C(2,0)\)时\(H=3(-2/3,4/3)-2(-1,1)=(-2,4)-(-2,2)=(0,2)\)；\(C(0,-2)\)时\(H=3(-4/3,2/3)-2(-1,1)=(-4,2)-(-2,2)=(-2,0)\)。（直接验证\(C(2,0)\)情形：\(AC\)水平\(\Longrightarrow\)过\(B\)的高为\(x=0\)；过\(C\)垂直\(AB\)的高\(y=-x+2\)；交点\((0,2)\)。）B对。\par\noindent 面积比：\(AB\)：\(x-y+4=0\)与欧拉线平行。\(C(2,0)\)到欧拉线距离\(4/\sqrt{2}\)，到\(AB\)距离\(6/\sqrt{2}\)，相似比\(2/3\)，小三角形与大三角形面积比\(4/9\)，故两部分之比\(4:5\)，即\(4/5\)，D对。选BCD。}
\end{ansblock}

\tihao{11}\zu{两直线位置关系的判定×填空} \tieside{简单(知识点三)}已知三条直线\(l_1\)：\(2x+my+2=0\)（\(m\in R\)），\(l_2\)：\(2x+y-1=0\)，\(l_3\)：\(x+ny+1=0\)（\(n\in R\)），若\(l_1\parallel l_2\)，\(l_1\perp l_3\)，则\(m+n\)的值为\kongda{−1}．

\begin{ansblock}[2章-练-课时06-T11]
% ans:2章-练-课时06-T11
\ansitem{11}{−1}
\ansnote{详解}{\(l_1\parallel l_2\Longleftrightarrow 2\cdot 1-2\cdot m=0\Longrightarrow m=1\)（此时\(l_1\)：\(2x+y+2=0\)与\(l_2\)不重合）。\(l_1\perp l_3\Longleftrightarrow 2\cdot 1+1\cdot n=0\Longrightarrow n=-2\)。\(m+n=-1\)。}
\end{ansblock}

\tihao{12}\zu{两直线位置关系的判定×填空} \tieside{中档(知识点二)}已知直线\(l_1\)：\(x+3y-5=0\)，\(l_2\)：\(3kx-y+1=0\)．若\(l_1\)，\(l_2\)与两坐标轴围成的四边形有一个外接圆，则\(k=\)\kongda{1}．

\begin{ansblock}[2章-练-课时06-T12]
% ans:2章-练-课时06-T12
\ansitem{12}{1}
\ansnote{详解}{四边形由\(l_1\)、\(l_2\)、\(x\)轴、\(y\)轴围成。坐标轴交点处的内角为\(90^\circ\)，四边形有外接圆\(\Longleftrightarrow\)对角互补\(\Longleftrightarrow l_1\)与\(l_2\)的交角为\(90^\circ\)，即\(l_1\perp l_2\)。\(k_1=-1/3\)，\(k_2=3k\)，垂直\(\Longleftrightarrow -k=-1\Longrightarrow k=1\)。（检验：\(l_1\)不过原点、不平行于坐标轴；\(k=1\)时\(l_2\)：\(3x-y+1=0\)同样，四边形存在且对角互补，有外接圆。）}
\end{ansblock}

\tihao{13}\zu{直线方程形式与互化×填空} \tieside{中档(知识点三)}（提示：本题第②空距离最大值用到§2.2.4点到直线的距离公式（06B 课时内容，位于本课时之后），供已提前选学该内容的学生使用；未学时可先做第①空，第②空待学完 06B 后回做。）已知点\(P(-2,0)\)和直线\(l\)：\((1+3\lambda)x+(1+2\lambda)y-(2+5\lambda)=0\)（\(\lambda\in R\)），该直线\(l\)过定点\kongda{(1,1)}，点\(P\)到直线\(l\)的距离\(d\)的最大值为\kongda{√10}．

\begin{ansblock}[2章-练-课时06-T13]
% ans:2章-练-课时06-T13
\ansitem{13}{(1,1)；√10}
\ansnote{详解}{按\(\lambda\)整理：\((x+y-2)+\lambda(3x+2y-5)=0\)，令两式同零：\(x+y-2=0\)且\(3x+2y-5=0\)，解得\(x=y=1\)，定点\(Q(1,1)\)。\(\lambda\)变化时\(l\)绕\(Q\)转动（斜率可取除\(-3/2\)外的一切值，且\(\lambda=-1/2\)时为铅直线\(x=1\)，故方向覆盖齐全）。\(d\)最大\(\Longleftrightarrow l\perp PQ\)，max\;\(d=|PQ|=\sqrt{(-2-1)^{2}+(0-1)^{2}}=\sqrt{10}\)。}
\end{ansblock}

\tihao{14}\zu{两直线位置关系的判定×填空} \tieside{中档(知识点三)}已知直线\(l_1\)经过点\(A(3,m)\)，\(B(m-1,2)\)，直线\(l_2\)经过点\(C(1,2)\)，\(D(-2,m+2)\)．
(1) 若\(l_1\parallel l_2\)，求实数\(m\)的值\kongda{m=1 或 6}；
(2) 若\(l_1\perp l_2\)，求实数\(m\)的值\kongda{m=3 或 −4}．

\begin{ansblock}[2章-练-课时06-T14]
% ans:2章-练-课时06-T14
\ansitem{14}{(1) m=1 或 6；(2) m=3 或 −4}
\ansnote{详解}{\(k_2=(m+2-2)/(-2-1)=-m/3\)；\(k_1=(2-m)/(m-1-3)=(2-m)/(m-4)\)。\par\noindent (1)\(k_1=k_2\)：\((2-m)/(m-4)=-m/3\Longrightarrow 3(2-m)=-m(m-4)\Longrightarrow m^{2}-7m+6=0\Longrightarrow m=1\)或6。检验不重合：\(m=1\)时\(l_1\)过\((3,1)\)、\((0,2)\)斜率\(-1/3\)，\(l_2\)过\((1,2)\)、\((-2,3)\)，\(C\)不在\(l_1\)上；\(m=6\)时\(l_1\)：\(y=-2x+12\)，\(C(1,2)\)不在其上。故\(m=1\)或6。\par\noindent (2)\(m=0\)时\(k_2=0\)（水平），需\(l_1\)铅直，但\(A(3,0)\)、\(B(-1,2)\)横坐标不同，斜率存在，舍；\(k_2\neq 0\)时\(k_1\cdot k_2=-1\)：\((2-m)/(m-4)\cdot(-m/3)=-1\Longrightarrow m(2-m)=-3(m-4)\Longrightarrow m^{2}-m-12=0\Longrightarrow m=3\)或\(-4\)。\(m=3\)：\(k_1=1\)，\(k_2=-1\)；\(m=-4\)：\(k_1=-3/4\)，\(k_2=4/3\)，积为\(-1\)。故\(m=3\)或\(-4\)。}
\end{ansblock}

% ============================================================
% 思维探索（题15~16＝T15~T16：解答2，居末档照库序；两问均 \liubai[24mm]）
% ============================================================
\dangbadge{思}{维}{探}{索}

\tihao{15}\zu{两直线位置关系的判定×解答} \tieside{中档(知识点三)}已知在\pxparallelogram\(ABCD\)中，\(A(1,2)\)，\(B(5,0)\)，\(C(3,4)\)．
(1) 求点\(D\)的坐标；
(2) 试判定\pxparallelogram\(ABCD\)是否为菱形？
\liubai[24mm]{解答书写区}

\begin{ansblock}[2章-练-课时06-T15]
% ans:2章-练-课时06-T15
\ansitem{15}{(1) D(−1,6)；(2) 是菱形}
\ansnote{详解}{(1) 平行四边形对角线互相平分：\(AC\)中点\(=BD\)中点。设\(D(x,y)\)：\(((1+3)/2,(2+4)/2)=((5+x)/2,(0+y)/2)\Longrightarrow x=-1\)，\(y=6\)，\(D(-1,6)\)。\par\noindent (2) 邻边相等即菱形：\(|AB|^{2}=16+4=20\)，\(|BC|^{2}=4+16=20\)，\(|AB|=|BC|\)，故为菱形。（或对角线：\(k_{AC}=1\)，\(k_{BD}=6/(-6)=-1\)，\(AC\perp BD\)且平分\(\Longrightarrow\)菱形。）}
\end{ansblock}

\tihao{16}\zu{直线方程形式与互化＋距离应用×解答} \tieside{中档(知识点三)}（提示：本题第(2)问求高（点到直线的距离）用到§2.2.4点到直线的距离公式（06B 课时内容，位于本课时之后），供已提前选学该内容的学生使用；未学时可用「等积法」（\(S=|BC|\cdot h/2=|x_By_C-x_Cy_B|/2\) 的向量/坐标展开）求面积，或待学完 06B 后回做。）已知\(\triangle ABC\)的三个顶点的坐标是\(A(1,1)\)，\(B(2,3)\)，\(C(3,-2)\)．
(1) 求\(BC\)边所在直线的方程；
(2) 求\(\triangle ABC\)的面积．
\liubai[24mm]{解答书写区}

\begin{ansblock}[2章-练-课时06-T16]
% ans:2章-练-课时06-T16
\ansitem{16}{(1) 5x+y−13=0；(2) 7/2}
\ansnote{详解}{(1) \(k_{BC}=(-2-3)/(3-2)=-5\)，\(y-3=-5(x-2)\)，即\(5x+y-13=0\)。\par\noindent (2) \(|BC|=\sqrt{1+25}=\sqrt{26}\)。\(A\)到\(BC\)的距离\(d=|5\cdot 1+1-13|/\sqrt{26}=7/\sqrt{26}\)。\(S=(1/2)\cdot\sqrt{26}\cdot(7/\sqrt{26})=7/2\)。}
\end{ansblock}

% ---- 尾块（\tailfill 置 \end{multicols} 前；无参＝内容尾随框，每件恰 1 处） ----
\tailfill

\end{multicols}
\end{document}
